\documentclass[10pt,letterpaper]{article}

\usepackage{fontspec}
\setmainfont{Latin Modern Roman}
\usepackage{amsmath,amssymb}
\usepackage[margin=0.8in]{geometry}
\usepackage{booktabs,longtable,array}
\usepackage{listings}
\usepackage{enumitem}
\usepackage[hidelinks]{hyperref}

\setlength{\parindent}{0pt}
\setlength{\parskip}{5pt}
\setlength{\emergencystretch}{3em}
\setlist[itemize]{leftmargin=1.5em,itemsep=2pt,topsep=4pt}
\setlist[enumerate]{leftmargin=1.7em,itemsep=3pt,topsep=4pt}
\renewcommand{\arraystretch}{1.1}

\providecommand{\tightlist}{%
  \setlength{\itemsep}{0pt}\setlength{\parskip}{0pt}}
\newcommand{\passthrough}[1]{#1}

\lstset{
  language=R,
  basicstyle=\ttfamily\small,
  frame=single,
  breaklines=true,
  showstringspaces=false,
  columns=fullflexible,
  keepspaces=true
}

\title{STAT 1210: Assignment 2\\[4pt]\large Part II: Probability and Distributions}
\author{Fall 2026 Semester}
\date{}

\begin{document}
\maketitle
\section*{Instructions}
\begin{itemize}
\tightlist
\item Your work should show all steps clearly for each question.
\item Your solutions must be handwritten on lined or letter sized paper. Your name and student number must appear on the first page.
\item Your name should also appear on every additional page and all pages should be stapled together.
\item R is the preferred programming language, although equivalent Python code is also acceptable unless a question specifically requires R code.
\item Programming resources, textbooks and software manuals may be consulted for coding tasks but all submitted work must be your own. Solutions generated by generative AI tools are not permitted.
\end{itemize}

\newpage

% =======================================
% Question 1
% =======================================
\subsection*{Question 1}
Let \(X\) be the number of seedlings, out of four, that show measurable root growth after a laboratory treatment. The proposed probability model is shown below.
\begin{longtable}[]{@{}lrrrrr@{}}
\toprule\noalign{}
	\(x\) & 0 & 1 & 2 & 3 & 4 \\
	\(P(X=x)\) & 0.18 & 0.31 & 0.29 & 0.16 & \(c\) 
\end{longtable}

\begin{enumerate}
\def\labelenumi{\arabic{enumi}.}
\tightlist
	\item What value must \(c\) have for this to be a valid probability model and which probability rule determines that value?
	\item What are \(P(X\ge 2)\), \(P(X<3)\), and \(P(X\ne 0)\)?
	\item What are the mean \(\mu_X\), variance \(\sigma_X^2\), and standard deviation \(\sigma_X\) of \(X\)? The formulas below may be used.
	\[
		\mu_X=\sum xP(X=x) \qquad \sigma_X^2=\sum (x-\mu_X)^2P(X=x)
	\]
	\item What does the value of \(\mu_X\) mean in the context of many repetitions of this four-seed experiment? Why does a non-integer expected value remain meaningful even though \(X\) can take only whole-number values?
	\item Verify your results using the following R code. Write the outputs of your code upto 6 decimal points. 
\begin{lstlisting}[language=R]
> x <- 0:4
> p <- c(0.18, 0.31, 0.29, 0.16, 0.06)
> sum(p)
> mu <- sum(x * p)
> variance <- sum((x - mu)^2 * p)
> c(mean = mu, variance = variance, sd = sqrt(variance))
\end{lstlisting}

\end{enumerate}

% =======================================
% Question 2
% =======================================
\subsection*{Question 2} 
The time \(T\), in minutes, required for an automated instrument to complete a calibration is modelled as uniformly distributed between 2 and 8 minutes. The density is therefore constant over the interval
\(2\le T\le8\).
\begin{enumerate}
\def\labelenumi{\arabic{enumi}.}
\tightlist
	\item What height must the uniform density curve have so that its total area equals 1?
	\item What are \(P(3.5\le T\le6.2)\) and \(P(T<2.8\text{ or }T>7.4)\)?
	\item What is \(P(T=5)\) and why does this answer not mean that a calibration time of exactly 5 minutes is impossible in practice?
	\item If the model is applied to 1200 independent calibrations, approximately how many calibration times would be expected to fall between 3.5 and 6.2 minutes? How does the Law of Large Numbers support this interpretation?
\end{enumerate}

\newpage

% =======================================
% Question 3
% =======================================
\subsection*{Question 3}
A monitoring program classified 400 water samples by region and by whether the contaminant concentration exceeded a review threshold.

\begin{longtable}[]{@{}lrrr@{}}
\toprule\noalign{}
Region & Below threshold & Above threshold & Total \\
\midrule\noalign{}
\endhead
\bottomrule\noalign{}
\endlastfoot
North & 84 & 16 & 100 \\
Central & 126 & 24 & 150 \\
South & 90 & 60 & 150 \\
\textbf{Total} & \textbf{300} & \textbf{100} & \textbf{400} \\
\end{longtable}

Let \(A\) denote ``the sample is above the threshold'' and \(S\) denote ``the sample is from the South region.''

\begin{enumerate}
\def\labelenumi{\arabic{enumi}.}
\tightlist
	\item What are the marginal probabilities \(P(A)\) and \(P(S)\)? What is the joint probability \(P(A\cap S)\)?
	\item What are \(P(A\mid S)\) and \(P(S\mid A)\)? How do their denominators and interpretations differ?
	\item What is the probability that a sample is either from the South region or above the threshold, \(P(S\cup A)\)?
	\item What are the conditional above-threshold rates within the North, Central and South regions? What evidence do these rates provide about an association between region and threshold status in this dataset?
	\item Why would these observational data alone be insufficient to conclude that region causes the difference in threshold status? What is one plausible additional variable that could be relevant?
\end{enumerate}

% =======================================
% Question 4
% =======================================
\subsection*{Question 4}
In a population of bacterial isolates, let \(A\) denote resistance to antibiotic A and \(B\) denote the presence of a particular plasmid marker. Suppose
\[
	P(A)=0.32 \qquad P(B)=0.25 \qquad P(A\cap B)=0.12
\]

\begin{enumerate}
\def\labelenumi{\arabic{enumi}.}
\tightlist
\item Are \(A\) and \(B\) disjoint? What feature of the given information determines the answer?
\item What is \(P(A\cup B)\)?
\item What is the probability that an isolate has neither antibiotic-A resistance nor the plasmid marker?
\item What are \(P(A\mid B)\) and \(P(B\mid A)\), and how should each probability be interpreted?
\item Are \(A\) and \(B\) independent? 
\end{enumerate}

% =======================================
% Question 5
% =======================================
\subsection*{Question 5}
A laboratory obtains sterile containers from two suppliers. Supplier A provides 65\% of the containers and has a 3\% defect rate. Supplier B provides 35\% of the containers and has an 8\% defect rate. One container is selected at random from the combined inventory.

\begin{enumerate}
\def\labelenumi{\arabic{enumi}.}
\tightlist
	\item How should a complete probability tree for supplier and defect status be drawn and labelled? The outgoing branches from every node should sum to 1.
	\item What are the four joint probabilities at the ends of the tree?
	\item What is the overall probability that the selected container is defective?
	\item Given that the selected container is defective, what is the probability that it came from Supplier B?
	\item Given that the selected container is not defective, what is the probability that it came from Supplier A?
	\item Are supplier and defect status independent under this model? What comparison supports the conclusion?
\end{enumerate}

\newpage

% =======================================
% Question 6
% =======================================
\subsection*{Question 6}
A screening test is used to detect a plant pathogen. In the population being screened, 6\% of plants carry the pathogen. The test has sensitivity 92\% and specificity 95\%. Assume these values apply to the population under consideration.

\begin{enumerate}
\def\labelenumi{\arabic{enumi}.}
\tightlist
	\item What do the sensitivity and specificity statements mean using conditional probability notation? Describe each in one or two sentences.
	\item In a hypothetical group of 10,000 plants, how many would be expected in each of the following categories: true positive, false negative, false positive and true negative?
	\item What is the overall probability that a randomly selected plant receives a positive result?
	\item Given a positive result, what is the probability that the plant carries the pathogen?
	\item Given a negative result, what is the probability that the plant carries the pathogen?
	\item Why is the sensitivity of 92\% not the answer to Part 4? How does the 6\% prevalence influence the interpretation of a positive result?
\end{enumerate}

% =======================================
% Question 6
% =======================================
\subsection*{Question 7}
Measurements from Laboratory A are approximately Normal with mean 120 units and standard deviation 15 units. Measurements from Laboratory B are approximately Normal with mean 80 units and standard deviation 5 units.

\begin{enumerate}
\def\labelenumi{\arabic{enumi}.}
\tightlist
	\item What are the \(z\)-scores of a value of 150 from Laboratory A and a value of 92 from Laboratory B?
	\item Which measurement is farther above the centre of its own distribution? What feature of the \(z\)-scores supports the answer?
	\item Under the Laboratory A model, what approximate proportion lies between 90 and 150 units according to the 68-95-99.7 rule?
	\item Under the Laboratory A model, what are the exact probabilities \(P(105<X<145)\) and \(P(X>150)\)? Use the following R code to calculate these probabilities and write the output upto to 6 decimal places as your solution. 
\begin{lstlisting}[language=R]
> pnorm(145, mean = 120, sd = 15) - pnorm(105, mean = 120, sd = 15)
> 1 - pnorm(150, mean = 120, sd = 15)
\end{lstlisting}

\end{enumerate}

% =======================================
% Question 8
% =======================================
\subsection*{Question 8} 
The mass \(M\) of a packaged laboratory standard is modelled as approximately Normal with mean 250 g and standard deviation 3 g.

\begin{enumerate}
\def\labelenumi{\arabic{enumi}.}
\tightlist
	\item What two mass values enclose the central 90\% of the model? The corresponding lower and upper cumulative probabilities should be identified before using \passthrough{\lstinline!qnorm()!}.
	\item What mass marks the beginning of the heaviest 1\% of packages?
	\item In a production run of 2,000 packages, approximately how many would be expected to exceed the cutoff from Part 2 if the model is accurate?
	\item How should each cutoff be interpreted in the original units and context?
	\item A histogram of recent package masses is strongly right-skewed. Why would the Normal-model cutoffs require caution even if the sample mean and standard deviation were close to 250 g and 3 g? Use the R code below and provide your solution upto 6 decimal places.
\begin{lstlisting}[language=R]
> qnorm(c(0.05, 0.95), mean = 250, sd = 3)
> qnorm(0.99, mean = 250, sd = 3)
\end{lstlisting}
\end{enumerate}

\newpage

% =======================================
% Question 9
% =======================================
\subsection*{Question 9}
For each situation below, which classification is most appropriate: \textbf{Binomial}, \textbf{Poisson} or \textbf{neither without additional assumptions}? For any situation classified as neither, what change or additional modelling condition might make a Binomial or Poisson approximation reasonable? Provide your reasons in 1-2 sentences.

\begin{enumerate}
\def\labelenumi{\arabic{enumi}.}
\tightlist
	\item A laboratory tests 20 independently selected samples. Each sample is classified as positive or negative, and the probability of a positive result is assumed to remain 0.08.
	\item A researcher counts bacterial colonies appearing on equal-area plates when colonies occur independently at a stable average rate per plate.
	\item Eight items are selected without replacement from a collection of 12 items containing five marked items, and the number of marked items selected is recorded.
	\item A service records the number of requests received each hour but the average request rate changes substantially between daytime and nighttime.
\end{enumerate}

% =======================================
% Question 10
% =======================================
\subsection*{Question 10}
An optical sensor passes a calibration check with probability 0.85. Twelve sensors are checked under a model in which outcomes are independent and the pass probability is the same for each sensor. Let \(X\) be the number that pass.

\begin{enumerate}
\def\labelenumi{\arabic{enumi}.}
\tightlist
	\item Why does \(X\) satisfy the four Binomial conditions, and how should its distribution be written?
	\item What are \(P(X=10)\) and \(P(X\ge10)\)? The first probability should be shown using the Binomial formula and both should be verified in R.
	\item What are the mean and standard deviation of \(X\)?
	\item How should the mean and standard deviation be interpreted across many repeated groups of 12 sensors?
	\item Why does an expected value of 10.2 not imply that 10.2 sensors can pass in one group?
\begin{lstlisting}[language=R]
> dbinom(10, size = 12, prob = 0.85)
> pbinom(9, size = 12, prob = 0.85, lower.tail = FALSE)
\end{lstlisting}
\end{enumerate}

% =======================================
% Question 11
% =======================================
\subsection*{Question 11}
Small surface defects occur along a manufactured sheet at an average rate of 2.4 defects per 100 metres. Assume a stable rate and independent counts in non-overlapping lengths. Let \(Y\) be the number of defects in a randomly selected 250-metre length.

\begin{enumerate}
\def\labelenumi{\arabic{enumi}.}
\tightlist
	\item Why is a Poisson model appropriate, and what value of \(\lambda\) applies to 250 metres?
	\item What are \(P(Y=4)\), \(P(Y\le2)\), and \(P(Y\ge1)\)? The probability notation should be translated carefully into \passthrough{\lstinline!dpois()!} or \passthrough{\lstinline ppois()!} commands.
	\item What are the mean, variance and standard deviation of \(Y\)?
	\item How would the value of \(\lambda\) change for a 50-metre length?
\end{enumerate}

% =======================================
% Question 12
% =======================================
\subsection*{Question 12}
A quality study models the number \(X\) of components requiring adjustment among 80 independently produced components as Binomial. The adjustment probability \(p\) is constant, \(p<0.50\) and the variance of \(X\) is 12.8.

\begin{enumerate}
\def\labelenumi{\arabic{enumi}.}
\tightlist
	\item Using \(\operatorname{Var}(X)=np(1-p)\), what quadratic equation must \(p\) satisfy?
	\item What are the two mathematical solutions for \(p\), and which solution is consistent with the information \(p<0.50\)?
	\item What are the mean and standard deviation of \(X\) under the selected model?
	\item What is \(P(14\le X\le18)\)? How can the inclusive interval be calculated using two cumulative Binomial probabilities?
	\item What is \(P(X\ge25)\), and would a count of 25 or more be unusual under this model? The numerical result and the model assumptions should both be considered.
\end{enumerate}


\end{document}